\chapter{Renormalization}

\begin{remarkbox}
Regularization makes divergent expressions well defined. Renormalization turns
regularized expressions into finite predictions by relating the parameters in the
Lagrangian to measured quantities. It is not just a trick for removing infinities;
it is the statement that physics depends on scale.
\end{remarkbox}

\section{The Basic Problem}

Perturbative loop calculations produce divergent expressions. After introducing a
regulator, these expressions depend on an auxiliary scale or cutoff. For example,
a one-loop correction may contain
\[
    \ln\frac{\Lambda}{m}
\]
with a cutoff regulator, or
\[
    \frac{1}{\epsilon}
\]
in dimensional regularization.

A physical prediction cannot depend on an arbitrary regulator. The way out is to
recognize that the parameters in the original Lagrangian are not directly the
measured parameters. Bare masses, fields, and couplings must be related to
renormalized quantities.

Renormalization is the systematic procedure for making this relation precise.

\section{Bare and Renormalized Quantities}

Start with a scalar theory. The bare Lagrangian may be written as
\[
    \mathcal{L}_0
    = \frac{1}{2}\partial_\mu\phi_0\partial^\mu\phi_0
      -\frac{1}{2}m_0^2\phi_0^2
      -\frac{\lambda_0}{4!}\phi_0^4.
\]
The subscript \(0\) denotes bare quantities. Introduce renormalized quantities by
\[
    \phi_0=Z_\phi^{1/2}\phi,
    \qquad
    m_0^2=m^2+\delta m^2,
    \qquad
    \lambda_0=\lambda+\delta\lambda
\]
with the precise form depending on convention.

The Lagrangian is rewritten as
\[
    \mathcal{L}_0=\mathcal{L}_{\mathrm{ren}}+\mathcal{L}_{\mathrm{ct}}.
\]
The renormalized part contains finite parameters \(m,\lambda\). The counterterm
part contains \(\delta Z_\phi,\delta m^2,\delta\lambda\), chosen to cancel loop
divergences.

\section{Counterterm Lagrangian}

For \(\phi^4\) theory, a common counterterm Lagrangian is
\[
    \mathcal{L}_{\mathrm{ct}}
    = \frac{1}{2}\delta Z_\phi\partial_\mu\phi\partial^\mu\phi
      -\frac{1}{2}\delta m^2\phi^2
      -\frac{\delta\lambda}{4!}\phi^4.
\]
Each counterterm gives an additional Feynman rule. For example:
\begin{itemize}
    \item \(\delta m^2\) gives a two-point insertion;
    \item \(\delta Z_\phi\) gives a momentum-dependent two-point insertion;
    \item \(\delta\lambda\) gives a four-point counterterm vertex.
\end{itemize}
The counterterm diagrams are included together with loop diagrams. Their role is
to cancel regulator-dependent divergent pieces.

This may look artificial at first, but it expresses a physical fact: the
parameters we write in a Lagrangian must be calibrated by measurements.

\section{Renormalization Conditions}

Counterterms are fixed by renormalization conditions. In an on-shell scheme for a
scalar field, one may require that the full propagator has a pole at the physical
mass:
\[
    p^2=m_{\mathrm{phys}}^2,
\]
and that the residue of the pole is one. These conditions define the physical
mass and the normalization of the field.

For the coupling, one may impose a condition on the four-point amplitude at a
chosen kinematic point. Schematically,
\[
    \mathcal{M}(s_0,t_0,u_0)=-\lambda_{\mathrm{ren}}.
\]
This equation defines the renormalized coupling at that reference point.

Different renormalization schemes impose different conditions. The numerical
values of renormalized parameters may differ from scheme to scheme, but physical
predictions agree when expressed consistently in terms of measured inputs.

\section{Two-Point Function and Self-Energy}

The full scalar propagator can be written as
\[
    G(p)=\frac{i}{p^2-m^2-\Pi(p^2)+i\epsilon},
\]
where \(\Pi(p^2)\) is the self-energy. The pole of the propagator is determined
by
\[
    p^2-m^2-\Pi(p^2)=0.
\]
The self-energy contains loop contributions and counterterm contributions.
Renormalization conditions fix the counterterms so that the pole occurs at the
chosen physical mass and with the chosen residue.

In perturbation theory, one expands \(\Pi\) order by order. At each order,
divergences are canceled by the counterterms allowed by the Lagrangian.

\section{Four-Point Function and Coupling Renormalization}

The four-point function determines the interaction strength. In \(\phi^4\)
theory, the tree-level amplitude is
\[
    \mathcal{M}_{\mathrm{tree}}=-\lambda.
\]
Loop corrections modify this amplitude. The renormalized four-point function has
schematic form
\[
    \mathcal{M}
    = -\lambda+\mathcal{M}_{\mathrm{loop}}+\mathcal{M}_{\mathrm{ct}}.
\]
The loop contribution is divergent before renormalization. The counterterm
contribution \(\mathcal{M}_{\mathrm{ct}}\) cancels the divergent part and fixes
how \(\lambda\) is related to a measured scattering amplitude.

This is coupling renormalization. It is the reason the coupling appearing in the
Lagrangian must be associated with a scale or a physical renormalization
condition.

\section{Field-Strength Renormalization}

The field-strength renormalization \(Z_\phi\) is fixed by the residue of the full
propagator. Near the physical pole,
\[
    G(p)\sim \frac{iZ}{p^2-m_{\mathrm{phys}}^2+i\epsilon}.
\]
In an on-shell scheme, one often chooses the renormalized field so that
\[
    Z=1.
\]
This condition determines \(\delta Z_\phi\).

Field-strength renormalization affects external legs in scattering amplitudes.
In LSZ reduction, the residue of the pole controls the normalization of
one-particle states. Thus \(Z_\phi\) is not merely a formal coefficient; it is
connected to how fields create physical particles from the vacuum.

\section{Minimal Subtraction}

In minimal subtraction, counterterms are chosen to remove only divergent pole
terms in dimensional regularization. For example, a counterterm may contain
\[
    \delta\lambda = \frac{a_1(\lambda)}{\epsilon}+\frac{a_2(\lambda)}{\epsilon^2}+\cdots.
\]
In the \(\overline{\mathrm{MS}}\) scheme, one also subtracts the standard
combination
\[
    -\gamma_E+\ln4\pi.
\]
Minimal subtraction schemes are not defined directly by physical pole masses or
scattering amplitudes. Instead, they define parameters at a renormalization scale
\(\mu\).

This makes them especially convenient for high-energy calculations and for the
renormalization group.

\section{Renormalization Scale}

Dimensional regularization introduces a scale \(\mu\). Renormalized parameters
therefore depend on \(\mu\):
\[
    m=m(\mu),
    \qquad
    \lambda=\lambda(\mu),
    \qquad
    e=e(\mu).
\]
Physical quantities cannot depend on the arbitrary scale \(\mu\). Therefore the
explicit \(\mu\)-dependence of loop corrections must cancel the implicit
\(\mu\)-dependence of renormalized parameters.

This requirement leads to renormalization group equations. For a coupling
\(g(\mu)\), one defines
\[
    \beta(g)=\mu\frac{\dd g}{\dd\mu}.
\]
The beta function tells us how the coupling changes with scale. The next chapter
will develop this idea systematically.

\section{Renormalizable and Nonrenormalizable Interactions}

A theory is perturbatively renormalizable if all divergences can be absorbed
into a finite number of counterterms already present in the Lagrangian. In four
dimensions, \(\phi^4\), Yukawa theory, and QED are renormalizable in this sense.

Interactions with couplings of negative mass dimension generate infinitely many
new counterterms at higher orders. Such theories are called nonrenormalizable in
the traditional sense. This does not mean they are useless. It means they should
be interpreted as effective field theories valid below some scale.

The modern viewpoint is that renormalizability is not simply a yes-or-no mark of
validity. It tells us how predictive a theory is at a given energy scale and how
many parameters are needed to reach a chosen accuracy.

\section{Physical Predictions}

A renormalized calculation follows a clear pattern:
\begin{enumerate}
    \item choose a regulator;
    \item compute loop diagrams and counterterm diagrams;
    \item impose renormalization conditions or choose a subtraction scheme;
    \item express bare parameters in terms of renormalized parameters;
    \item compute observables in terms of measured inputs.
\end{enumerate}
The final observable must be finite and independent of the regulator.

For example, a scattering cross section should not depend on \(1/\epsilon\) or on
a cutoff \(\Lambda\). It may depend on physical masses, physical charges,
external momenta, and the chosen renormalization scale in a way compensated by
running parameters.

\section{What Renormalization Means Physically}

Renormalization teaches that parameters are scale-dependent summaries of physics.
A mass or coupling in a Lagrangian is not an isolated number with meaning at all
scales. It is defined by a measurement procedure or a subtraction convention.

Short-distance fluctuations affect long-distance parameters. Instead of trying
to remove this fact, renormalization organizes it. The same theory can be
described at different scales with different parameter values while predicting
the same physical observables.

This is why renormalization is central to modern QFT. It is not merely a repair
of infinities; it is a theory of how physics changes with scale.

\begin{summarybox}
Renormalization relates bare Lagrangian parameters to finite measured quantities.
Loop divergences are canceled by counterterms, whose finite parts are fixed by a
renormalization scheme or physical conditions. Two-point functions determine mass
and field-strength renormalization, while four-point functions determine coupling
renormalization. The dependence of renormalized parameters on the scale \(\mu\)
leads to the renormalization group.
\end{summarybox}

\section{Chapter Checklist}

After reading this chapter, one should be able to:
\begin{itemize}
    \item distinguish bare and renormalized quantities;
    \item write a counterterm Lagrangian for \(\phi^4\) theory;
    \item explain how renormalization conditions fix counterterms;
    \item relate the two-point function to mass renormalization;
    \item relate the pole residue to field-strength renormalization;
    \item explain coupling renormalization through the four-point function;
    \item distinguish on-shell and minimal subtraction schemes;
    \item explain the role of the renormalization scale \(\mu\);
    \item define the beta function qualitatively;
    \item distinguish renormalizable interactions from effective interactions;
    \item explain why renormalization is about scale dependence, not only about
    removing infinities.
\end{itemize}
